\chapter{Introduction}\label{ch:introduction}
      \section{Overview}
        In this project, I will construct and analyse supersingular $\ell$-isogeny graphs over finite fields $\mathbb{F}_{p^2}$. These are graphs whose vertices represent isomorphism classes of supersingular elliptic curves over $\mathbb{F}_{p^2}$ and whose edges represent degree-$\ell$ isogenies between them (see \cite{isogeny_school1} for background). These graphs are of special interest in cryptography, since many cryptographic protocols can be understood as applying random walks on these graphs. One such example is given in \cite{CSIDH}, whose designs provide inspiration for the graphs considered in this project.

        The significance of these graphs in cryptography arises mainly from their strong expansion properties; in particular, they are Ramanujan graphs, and therefore random walks on these graphs converge quickly to a uniform distribution. The most cryptographically interesting consequence is that given any two vertices in one such graph, it will be computationally expensive to find a path connecting them.

        We are therefore interested in performing empirical verification of theoretical properties attributed to these graphs. To this end, two main interests stand out: the eigenvalue spectra associated to the adjacency matrices of these graphs, and the mixing times of random walks on these graphs. For this project, I will produce a sample of supersingular isogeny graphs over a range of parameters and compute these two characteristics, comparing them to theoretical estimates in order to illustrate and corroborate them.

      \section{Scope}
        We have the following primary objectives. 
        \begin{enumerate}
            \item 
                Efficiently compute supersingular $\ell$-isogeny graphs over $\mathbb{F}_{p^2}$ for a range of parameters $p$ and $\ell$.
            \item 
                Compute the adjacency matrices for a sample of supersingular isogeny graphs and compute their eigenvalue spectra. We will use this to answer two questions.
                \begin{enumerate}
                    \item 
                        How does the spectral gap of a supersingular $\ell$-isogeny graph over $\mathbb{F}_{p^2}$ vary according to the parameters $\ell$ and $p$?
                    \item 
                        How does the spectral gap of a supersingular $\ell$-isogeny graph over $\mathbb{F}_{p^2}$, as measured empirically, align with theoretical estimates of the same quantity?
                \end{enumerate}
            \item 
                Produce, for each sampled supersingular $\ell$-isogeny graph over $\mathbb{F}_{p^2}$, a sample of random walks on that graph. We will use this to answer two questions.
                \begin{enumerate}
                    \item
                        How are mixing times of random walks distributed for each graph, in terms of parameters $\ell$ and $p$?
                    \item 
                        How do the mixing times of random walks on these graphs, as measured empirically, align with theoretical estimates of mixing times for Ramanujan graphs?
                \end{enumerate}
        \end{enumerate}

    
    